\section[Spin weights of an antiferromagnet]{Linear spin-wave weights of a two-sublattice antiferromagnet}\label{app:D.38}

The notation is that of \cref{sec:7.1}: sublattices A and B of $N_{\mathrm{s}}$ sites each, Holstein--Primakoff bosons $d_{\mathrm{A}i}$ and $d_{\mathrm{B}j}$, the coefficients $A_{\vect{k}}$ and $B_{\vect{k}}$ of \cref{eq:7.1-Hk}, the magnon operators $a_{\vect{k}1}$ and $a_{\vect{k}2}$ of \cref{eq:7.1-bogo}, and the ground state $|0\rangle$ that they annihilate. We take B to be A translated by a fixed vector $\vect{\delta}_0$, as in NiO and in the simple-cubic antiferromagnet; then $2\vect{\delta}_0$ is a vector of sublattice A.

\subsection*{Transverse matrix elements}

A momentum transfer $\vect{q}$ outside the Brillouin zone of one sublattice is split as $\vect{q} = \vect{q}_r + \vect{G}$, with $\vect{q}_r$ in that zone and $\vect{G}$ a reciprocal-lattice vector of the sublattice, as in \cref{eq:3.1-qsplit}. The sites of A are lattice vectors of the sublattice, so $\ee^{\ii\vect{G}\cdot\vect{R}_i} = 1$ on A, while on B, $\ee^{\ii\vect{G}\cdot\vect{R}_j} = \ee^{\ii\vect{G}\cdot\vect{\delta}_0} \equiv \epsilon_{\vect{G}}$. Since $2\vect{\delta}_0$ lies on A, $\epsilon_{\vect{G}}^2 = 1$ and $\epsilon_{\vect{G}} = \pm1$. The same phase relates the lattice sums of \cref{eq:D.19-Jk} at $\vect{q}$ and at $\vect{q}_r$:
\begin{equation*}
  J^{\mathrm{ex}}_{\mathrm{AA}}(\vect{q}) = J^{\mathrm{ex}}_{\mathrm{AA}}(\vect{q}_r), \qquad J^{\mathrm{ex}}_{\mathrm{AB}}(\vect{q}) = \epsilon_{\vect{G}}\,J^{\mathrm{ex}}_{\mathrm{AB}}(\vect{q}_r) ,
\end{equation*}
because every neighbor on A differs from the origin by a vector of A, and every neighbor on B by $\vect{\delta}_0$ plus a vector of A. Defined by the same formulas at the full $\vect{q}$, the coefficients of \cref{eq:7.1-Hk} are $A_{\vect{q}} = A_{\vect{q}_r}$ and $B_{\vect{q}} = \epsilon_{\vect{G}}B_{\vect{q}_r}$.

At linear order $S^+_i = (2|\vect{S}|)^{1/2}d_{\mathrm{A}i}$ on A and $S^+_j = (2|\vect{S}|)^{1/2}d^\dagger_{\mathrm{B}j}$ on B (\cref{sec:7.1}). In the Fourier transform, the lattice sum over A picks $\vect{k} = -\vect{q}_r$ and the lattice sum over B carries the extra factor $\epsilon_{\vect{G}}$. The inverse transformation \cref{eq:D.19-inverse}, with $u^{\mathrm{M}}_{\vect{k}}$ and $v^{\mathrm{M}}_{\vect{k}}$ real and even in $\vect{k}$, then gives
\begin{equation*}
  \begin{aligned}
  S^+(\vect{q}) &= (2|\vect{S}|N_{\mathrm{s}})^{1/2}\bigl[(u^{\mathrm{M}} - \epsilon_{\vect{G}}v^{\mathrm{M}})\,a_{-\vect{q}_r1} + (\epsilon_{\vect{G}}u^{\mathrm{M}} - v^{\mathrm{M}})\,a^\dagger_{-\vect{q}_r2}\bigr] , \\
  S^-(\vect{q}) &= (2|\vect{S}|N_{\mathrm{s}})^{1/2}\bigl[(u^{\mathrm{M}} - \epsilon_{\vect{G}}v^{\mathrm{M}})\,a^\dagger_{\vect{q}_r1} + (\epsilon_{\vect{G}}u^{\mathrm{M}} - v^{\mathrm{M}})\,a_{\vect{q}_r2}\bigr] ,
  \end{aligned}
\end{equation*}
with the coefficients at $\vect{q}_r$. On the ground state, $S^-(\vect{q})$ creates the magnon $a^\dagger_{\vect{q}_r1}$ and $S^+(\vect{q})$ the magnon $a^\dagger_{-\vect{q}_r2}$, both of energy $\omega^{\mathrm{M}}_{\vect{q}_r}$, with amplitudes $c$ and $\epsilon_{\vect{G}}c$, where $c = (2|\vect{S}|N_{\mathrm{s}})^{1/2}(u^{\mathrm{M}}_{\vect{q}_r} - \epsilon_{\vect{G}}v^{\mathrm{M}}_{\vect{q}_r})$. The expressions of \cref{sec:7.1} for $(u^{\mathrm{M}})^2$, $(v^{\mathrm{M}})^2$ and $u^{\mathrm{M}}v^{\mathrm{M}}$ give
\begin{equation}
  w^{\mathrm{M}}_{\vect{q}} \equiv \bigl(u^{\mathrm{M}}_{\vect{q}_r} - \epsilon_{\vect{G}}v^{\mathrm{M}}_{\vect{q}_r}\bigr)^2
  = \frac{A_{\vect{q}_r} - \epsilon_{\vect{G}}B_{\vect{q}_r}}{(A_{\vect{q}_r}^2 - B_{\vect{q}_r}^2)^{1/2}}
  = \frac{|\vect{S}|\,(A_{\vect{q}} - B_{\vect{q}})}{\omega^{\mathrm{M}}_{\vect{q}}} ,
  \label{eq:D.38-weight}
\end{equation}
where the last form uses the coefficients at the full momentum transfer, which absorb the phase. For $\epsilon_{\vect{G}} = \pm1$, $w^{\mathrm{M}}_{\vect{q}}$ is the weight $(u^{\mathrm{M}} \mp v^{\mathrm{M}})^2$ of \cref{eq:7.1-weights}, with the upper sign near $\Gamma$ and the lower sign near the ordering wavevector.

The spin density is $\vect{s}_e(\vect{q}) = F^{\mathrm{mag}}(q)\sum_j\ee^{\ii\vect{q}\cdot\vect{R}_j}\vect{S}_j$ (\cref{eq:7.1-Fmag}), with $s_{e,x} = F^{\mathrm{mag}}(S^+ + S^-)/2$ and $s_{e,y} = F^{\mathrm{mag}}(S^+ - S^-)/2\ii$, and $S^z$ creates no single magnon. The two branches have the matrix elements $(\langle s_{e,x}\rangle, \langle s_{e,y}\rangle) = F^{\mathrm{mag}}c\,(1, \ii)/2$ and $F^{\mathrm{mag}}\epsilon_{\vect{G}}c\,(1, -\ii)/2$, which rotate in opposite senses. When these matrix elements are inserted in \cref{eq:4.2-Sab} at $T = 0$ with $\Vol = 2N_{\mathrm{s}}\Omega_{\mathrm{ion}}$, the off-diagonal elements of the two branches cancel, and with $|c|^2 = 2|\vect{S}|N_{\mathrm{s}}w^{\mathrm{M}}_{\vect{q}}$ the diagonal elements give \cref{eq:14.1-Safm}. With $\hat{\vect{z}}$ replaced by the local ordering axis, \cref{eq:14.1-Safm} holds for any orientation of the moments.

\subsection*{Transverse weight and zero-point reduction}

The equal-time weight of a response is its integral over $\omega$, and the equal-time weight summed over $\vect{q}$ in the zone becomes a correlation of one site. For the transverse spin the site correlation is direct: on A, $\frac12\langle S^+_iS^-_i + S^-_iS^+_i\rangle = |\vect{S}|\langle d_{\mathrm{A}i}d^\dagger_{\mathrm{A}i} + d^\dagger_{\mathrm{A}i}d_{\mathrm{A}i}\rangle$ at linear order. The deviation per site follows from \cref{eq:D.19-inverse}: $\langle d^\dagger_{\mathrm{A}\vect{k}}d_{\mathrm{A}\vect{k}}\rangle = (v^{\mathrm{M}}_{\vect{k}})^2$, since only the term $(v^{\mathrm{M}}_{\vect{k}})^2a_{\vect{k}2}a^\dagger_{\vect{k}2}$ survives in the vacuum. Hence
\begin{equation}
  \begin{gathered}
  \Delta S \equiv \langle d^\dagger_{\mathrm{A}i}d_{\mathrm{A}i}\rangle = \frac{1}{N_{\mathrm{s}}}\sum_{\vect{k}}\bigl(v^{\mathrm{M}}_{\vect{k}}\bigr)^2, \qquad
  \langle S^z_i\rangle = |\vect{S}| - \Delta S, \\
  W_\perp \equiv \langle S^x_iS^x_i + S^y_iS^y_i\rangle = |\vect{S}|\,(1 + 2\Delta S) ,
  \end{gathered}
  \label{eq:D.38-Wperp}
\end{equation}
the same on B. The zero-point reduction $\Delta S$ does not depend on $|\vect{S}|$, since $A_{\vect{k}}$ and $B_{\vect{k}}$ do not. The same $W_\perp$ follows from \cref{eq:D.38-weight}. As $\vect{q}$ runs over a full period of the reciprocal lattice of all the magnetic ions, A and B together, each $\vect{q}_r$ appears once with $\epsilon_{\vect{G}} = +1$ and once with $-1$. The average of $w^{\mathrm{M}}_{\vect{q}}$ over this period is therefore $(u^{\mathrm{M}})^2 + (v^{\mathrm{M}})^2 = 1 + 2(v^{\mathrm{M}})^2$, which averages to $1 + 2\Delta S$ over the zone. Integrating $S_{xx} + S_{yy}$ of \cref{eq:14.1-Safm} over $\omega/2\pi$, multiplying by $\Omega_{\mathrm{ion}}$ and averaging over $\vect{q}$ then gives $|\vect{S}|(1 + 2\Delta S)$ at $F^{\mathrm{mag}} = 1$. The one-magnon line carries all of $W_\perp$ at this order.

\subsection*{Longitudinal weight}

On A, $S^z_i - \langle S^z_i\rangle = -(n_i - \Delta S)$ with $n_i = d^\dagger_{\mathrm{A}i}d_{\mathrm{A}i}$. By \cref{eq:D.19-inverse}, $n_i$ contains terms $a^\dagger_1a^\dagger_2$, which acting on $|0\rangle$ create two magnons, and terms that annihilate it, apart from the constant $\Delta S$. The connected longitudinal weight is therefore the two-magnon weight. The ground state of a quadratic boson Hamiltonian is Gaussian, and Wick's theorem gives
\begin{equation*}
  \langle n_i^2\rangle = \langle d^\dagger d\rangle^2 + \langle d^\dagger d\rangle\langle dd^\dagger\rangle + \langle d^\dagger d^\dagger\rangle\langle dd\rangle = \Delta S^2 + \Delta S(1 + \Delta S) ,
\end{equation*}
since the anomalous average on one site vanishes: $\langle d_{\mathrm{A}\vect{k}}d_{\mathrm{A}\vect{k}'}\rangle = -u^{\mathrm{M}}_{\vect{k}}v^{\mathrm{M}}_{\vect{k}'}\langle a_{\vect{k}1}a^\dagger_{\vect{k}'2}\rangle = 0$, the two branches being independent. The pairing of \cref{eq:7.1-Hk} correlates a deviation on A only with a deviation on B. Hence
\begin{equation}
  W_\parallel \equiv \langle(S^z_i)^2\rangle - \langle S^z_i\rangle^2 = \Delta S\,(1 + \Delta S), \qquad
  \frac{W_\parallel}{W_\perp} = \frac{\Delta S\,(1 + \Delta S)}{|\vect{S}|\,(1 + 2\Delta S)} ,
  \label{eq:D.38-Wpar}
\end{equation}
which is \cref{eq:14.1-ratio}.

\subsection*{The spin sum rule}

The identity $\vect{S}_i^2 = |\vect{S}|(|\vect{S}| + 1)$ holds on every state, so the three weights and the ordered moment obey $W_\perp + W_\parallel + \langle S^z_i\rangle^2 = |\vect{S}|(|\vect{S}| + 1)$ exactly. With \cref{eq:D.38-Wperp,eq:D.38-Wpar} inserted, the left side is $|\vect{S}|(|\vect{S}| + 1) + \Delta S(1 + 2\Delta S)$: linear spin-wave theory overcounts the total weight by $\Delta S(1 + 2\Delta S)$ per site. With $\langle S^z_i\rangle$ and $W_\parallel$ as above, the identity requires
\begin{equation}
  W_\perp^{\mathrm{sr}} = |\vect{S}|(|\vect{S}| + 1) - (|\vect{S}| - \Delta S)^2 - \Delta S(1 + \Delta S) = (|\vect{S}| - \Delta S)(1 + 2\Delta S) ,
  \label{eq:D.38-sumrule}
\end{equation}
a fraction $1 - \Delta S/|\vect{S}|$ of the linear spin-wave value. The excess weight $\Delta S(1 + 2\Delta S)$ comes from the square roots of \cref{eq:7.1-HP}, which linear spin-wave theory replaces by $(2|\vect{S}|)^{1/2}$. Kept whole, the square roots give $\langle S^+_iS^-_i\rangle = \langle(2|\vect{S}| - n_i)(n_i + 1)\rangle$ and $\langle S^-_iS^+_i\rangle = \langle n_i(2|\vect{S}| + 1 - n_i)\rangle$; the half-sum of these two expectation values, evaluated in the same Gaussian state with $\langle n_i\rangle = \Delta S$ and $\langle n_i^2\rangle = \Delta S(1 + 2\Delta S)$, returns \cref{eq:D.38-sumrule}. The first omitted term of the transverse operator, $-(2|\vect{S}|)^{1/2}n_id_{\mathrm{A}i}/4|\vect{S}|$, removes the excess to this order. It also creates three magnons, and the weight of these three-magnon states lowers the one-magnon part of $W_\perp^{\mathrm{sr}}$ further.
